\section{Discussion}
\label{sec:discussion}

% =============================================================================
% 5.1  Physical interpretation: why frequency drops with ripening
% =============================================================================
\subsection{Why the tap-tone drops with ripening}
\label{sec:physical_interp}

The monotonic decrease in tap-tone frequency during ripening
has a straightforward physical
explanation.  As the fruit ripens, enzymatic degradation of pectin in the
rind cell walls reduces the effective Young's modulus
$E_{\mathrm{rind}}$.  Since $f_2 \propto \sqrt{E_{\mathrm{rind}}}$
(Eq.~\ref{eq:f2_squared}), the frequency falls.  For the canonical
Crimson Sweet geometry, $f_2$ drops from \SI{203}{\hertz} (unripe,
$E = \SI{200}{\mega\pascal}$) to \SI{46}{\hertz} (overripe,
$E = \SI{15}{\mega\pascal}$)---a factor of~4.4 reduction.  This is
precisely the information that experienced shoppers extract
unconsciously: a lower pitch signals a softer, riper melon.

The frequency shift is not driven by changes in fluid density or fruit
geometry, both of which vary by less than $\SI{5}{\percent}$ over the
ripening window.  The Sobol analysis (\S\ref{sec:sobol}) confirms that
$E_{\mathrm{rind}}$ is the dominant parameter, contributing
\SI{54}{\percent} of the total variance ($S_T = 0.54$), whereas flesh
density and turgor pressure contribute less than \SI{0.1}{\percent}
each.  The secondary contribution of fruit size ($S_T = 0.38$) is
relevant across cultivars but approximately constant during ripening of
a single fruit.

The quality factor also decreases during ripening---from $Q = 12.5$
(unripe) to $Q = 5.0$ (overripe)---as the loss tangent increases from
0.08 to~0.20 with pectin degradation.  In practice, this means that an
overripe melon produces a duller, more heavily damped thump, consistent
with the subjective experience of a `mushy' sound.

% =============================================================================
% 5.2  Practical workflow and measurement conditions
% =============================================================================
\subsection{Practical workflow and measurement considerations}
\label{sec:practical_workflow}

Figure~\ref{fig:workflow} illustrates the envisaged measurement
workflow.  The procedure involves five steps: (1)~tap the fruit with a
knuckle or small impactor; (2)~record the acoustic response with a
microphone (or accelerometer attached to the rind); (3)~extract the
dominant spectral peak~$f_2$ via FFT; (4)~input the fruit geometry
($a$, $c$, $h$), obtainable from callipers or a photograph; and
(5)~apply the closed-form inversion (Eq.~\ref{eq:inversion}) to
estimate~$E_{\mathrm{rind}}$.

\begin{figure}[t]
  \centering
  \fbox{\parbox{0.85\textwidth}{\centering\small
    \textbf{Step 1: Tap} $\longrightarrow$
    \textbf{Step 2: Record} $\longrightarrow$
    \textbf{Step 3: Extract $f_2$} $\longrightarrow$
    \textbf{Step 4: Input geometry} $\longrightarrow$
    \textbf{Step 5: Invert $\to E_{\mathrm{rind}}$}\\[6pt]
    Knuckle / impactor \quad$\to$\quad
    Microphone / accelerometer \quad$\to$\quad
    FFT peak detection \quad$\to$\quad
    $(a, c, h)$ from callipers or photo \quad$\to$\quad
    Eq.~\eqref{eq:inversion}
  }}
  \caption{Envisaged practical workflow for acoustic rind stiffness
    estimation.  The closed-form inversion requires only the
    tap-tone frequency~$f_2$ and approximate fruit geometry.
    The subsequent mapping from $E_{\mathrm{rind}}$ to a ripeness
    classification (dashed) requires cultivar-specific calibration
    data not developed in this study.}
  \label{fig:workflow}
\end{figure}

Several practical measurement conditions merit discussion:
\begin{itemize}
  \item \textbf{Support conditions.}  The fruit should rest on a soft
    support (e.g.\ foam pad) to approximate free boundary conditions.
    A hard surface introduces contact stiffness that may shift or
    split the $n = 2$ resonance.
  \item \textbf{Tap location.}  The $n = 2$ flexural mode has
    maximum displacement at the equator; tapping near the equatorial
    plane maximises the modal excitation.  Tapping at the poles
    preferentially excites axisymmetric ($n = 0$) breathing modes at
    much higher frequencies.
  \item \textbf{Microphone or sensor position.}  The microphone
    should be placed near the fruit surface (within a few centimetres)
    to maximise the signal-to-noise ratio.  An accelerometer bonded
    to the rind provides a stronger signal but requires contact.
  \item \textbf{Ambient noise.}  Outdoor or market environments may
    mask the tap-tone signal; a brief time-windowed acquisition
    ($\sim\SI{50}{\milli\second}$) centred on the tap event can
    reject most background noise.
\end{itemize}

% =============================================================================
% 5.3  Limitations
% =============================================================================
\subsection{Limitations}
\label{sec:limitations}

Several simplifications deserve acknowledgement:

\begin{itemize}
  \item \textbf{Equivalent-sphere approximation.}  The present model
    captures geometry through the equivalent-volume sphere
    radius~$R_{\mathrm{eq}} = (a^2 c)^{1/3}$ and does not include
    explicit curvature corrections for oblate deformation.  All
    stiffness and inertia expressions
    (Eqs.~\ref{eq:K_bend}--\ref{eq:K_prestress},~\ref{eq:rho_eff})
    are those of a spherical shell.  Two watermelons with the same
    $R_{\mathrm{eq}}$ but different aspect ratios~$\zeta$ are therefore
    assigned identical frequencies.  For watermelons with moderate
    aspect ratios ($\zeta \approx 0.7$--$0.9$), this approximation
    introduces an error of unquantified magnitude in the predicted
    tap-tone frequency.  The dominant physics is expected to be
    controlled by the shell's average radius rather than its local
    curvature, but bounding this error rigorously would require
    comparison to a Ritz variational model
    with explicit $\zeta$-dependent basis functions (e.g.\ Legendre
    polynomials in oblate spheroidal coordinates), which is identified
    as future work.

  \item \textbf{Isotropic rind.}  The watermelon rind is a biological
    composite with fibre orientation that may induce anisotropy.  We
    treat it as isotropic; orthotropic extensions are straightforward
    but require additional material characterisation.

  \item \textbf{Uniform thickness.}  Real rinds are thicker near the
    blossom end.  Thickness variation could be incorporated via
    perturbation methods but is unlikely to shift $f_2$ by more than
    $\sim\SI{10}{\percent}$, based on the $\pm\SI{10}{\percent}$
    sensitivity analysis in \S\ref{sec:geom_uncertainty}
    ($\Delta f_2 \in [\SI{-4.6}{\percent}, \SI{+4.4}{\percent}]$ for
    $\pm\SI{10}{\percent}$ uniform thickness perturbation).

  \item \textbf{Flesh as inviscid fluid.}  Watermelon flesh has finite
    shear strength.  At the frequencies of interest
    ($\sim\SI{100}{\hertz}$), the flesh is well within the
    fluid-like regime ($\omega \tau \ll 1$, where $\tau$ is the
    viscoelastic relaxation time), so this approximation is reasonable.

  \item \textbf{Turgor pressure.}  We neglect internal turgor pressure
    effects on the eigenfrequency in the inversion.  The Sobol analysis
    confirms that $P_{\mathrm{int}}$ contributes less than
    \SI{0.1}{\percent} of the total variance in $f_2$, validating this
    simplification.  For very unripe fruit with intact cell walls,
    turgor may contribute a correction of order $\SI{10}{\percent}$,
    but this falls within the geometric uncertainty band.

  \item \textbf{Experimental validation.}  The model predictions are
    compared with published data (\S\ref{sec:published_data}) but have
    not yet been validated against purpose-designed experiments with
    controlled geometry and material characterisation.  Such validation
    is a priority for future work.

  \item \textbf{Frequency-to-ripeness chain.}  The model establishes
    the link from tap-tone frequency to effective rind modulus.  The
    further link from modulus to biological ripeness (eating quality,
    sugar content, texture) requires cultivar-specific calibration
    against destructive reference measurements, which has not been
    performed here.  The model therefore delivers a \emph{stiffness
    parameter}, not a ripeness classification.
\end{itemize}

% =============================================================================
% 5.4  Practical implications
% =============================================================================
\subsection{Practical implications}
\label{sec:practical}

The closed-form inversion (Eq.~\ref{eq:inversion}) is computationally
trivial and could be implemented in a smartphone application.  A
microphone records the tap-tone, a fast Fourier transform extracts
$f_2$, and the inversion formula returns $E_{\mathrm{rind}}$ given
approximate fruit dimensions (obtainable from a photograph).  The
dimensionless geometric invariant $\Pi_{\mathrm{ripe}}$ then provides
an approximate cultivar-independent calibration constant that maps the
result onto a modulus scale, although the \SI{5.4}{\percent} spread
across cultivars (Table~\ref{tab:Pi_ripe}) means that cultivar-specific
calibration would be preferable in practice.  We reiterate that
$\Pi_{\mathrm{ripe}}$ collapses geometries, not ripeness stages;
its utility as a ripeness predictor has not been biologically validated.

The Monte Carlo uncertainty analysis (\S\ref{sec:geom_uncertainty})
demonstrates that even with $\pm\SI{10}{\percent}$ geometric
measurement error, the inverted modulus has a coefficient of variation
of only \SI{13.5}{\percent} ($\SI{95}{\percent}$ CI:
$[\SI{38.5}{\mega\pascal}, \SI{64.1}{\mega\pascal}]$ for a true value
of \SI{50}{\mega\pascal}).  This uncertainty is well within the
factor-of-four modulus difference between ripeness categories
(Table~\ref{tab:forward_results}), suggesting that reliable
\emph{stiffness discrimination} across ripeness categories is
achievable even with imprecise geometry.  (Whether this stiffness
discrimination translates to reliable ripeness classification
requires cultivar-specific calibration.)

This approach offers several advantages over existing non-destructive
methods (near-infrared spectroscopy, impact testing with accelerometers):
it requires no specialised hardware, is non-contact (or nearly so---one
need only tap), and provides a physics-based rather than purely
statistical estimate of internal quality.  The Sobol analysis further
shows that the parameters most difficult to measure \emph{in situ}
(flesh density, Poisson's ratio, turgor pressure) are precisely those
to which the tap-tone is least sensitive, a fortuitous alignment that
makes the method practically robust.

% =============================================================================
% 5.5  Generality of the framework
% =============================================================================
\subsection{Generality of the shell--fluid framework}
\label{sec:generality}

The model is not fruit-specific.  Any fluid-filled biological structure
that can be approximated as a thin shell falls within its
scope.  Candidate applications include other cucurbits (cantaloupe,
honeydew, pumpkin), drupes (coconut), and even non-biological systems
such as pressurised tanks and sporting goods.  The multi-cultivar
comparison (\S\ref{sec:cultivar}) demonstrates that cultivar-specific
geometry is easily accommodated by adjusting three parameters ($a$, $c$,
$h$), and the dimensionless geometric invariant~$\Pi_{\mathrm{ripe}}$
provides a natural framework for cross-species calibration.

The collapse of $\Pi_{\mathrm{ripe}}$ to within
\SI{3.3}{\percent} across four geometrically distinct watermelon
cultivars (Table~\ref{tab:Pi_ripe}) suggests that extension to other
fruit species requires only measurement of the species-specific
shell geometry ($a$, $c$, $h$), which determines~$R_{\mathrm{eq}}$
and thereby the frequency scaling.

% =============================================================================
% 5.6  Connection to biomedical vibroacoustics
% =============================================================================

\paragraph{Connection to biomedical vibroacoustics.}
The equivalent-sphere shell framework used here is structurally
identical to the model applied to the human abdomen in our companion
work~\cite{MaceAndMace2026}, where the abdominal wall
($E \sim \SI{0.1}{\mega\pascal}$) encloses a fluid-filled cavity with
eigenfrequencies near \SI{4}{\hertz}.  The watermelon rind
($E \sim \SIrange{15}{200}{\mega\pascal}$) shifts these frequencies
upward to \SIrange{46}{203}{\hertz}, but the governing equations are
unchanged.  This parallel illustrates that classical shell
theory~\cite{lamb1882, Junger1986} provides a common analytical basis
for biological fluid-filled structures across very different material
scales, from clinical diagnostics to agricultural quality assessment.
